\scanpage{216}
\mspage{132}
% archive sheet number: 215
% element: text | top, lines 1-8 | Prouvons 2 degrees
\uline{Prouvons 2\textsuperscript{o}}. La première \deleted{\unsure{implication}} inclusion (4.5.13.1) est tautologique, la deuxième ne fait qu'exprimer $P(\pi)\Rightarrow d$. Le fait que $\pi^{\pi_0}=\pi$ implique l'égalité des conditions extrêmes est un rappel de 4.5.6. Reste à prouver que $\mathrm{Filt}'{}^\pi$ est stable dans $M$ par sommes de cardinal $\leq\pi$. Mais
% element: text | middle, lines 9-17 | si x alpha
si $(x_\alpha)_{\alpha\in A}$ est une famille d'objets de $\mathrm{Filt}'{}^\pi$, avec $\operatorname{card}A\leq\pi$, et si $x=\coprod_{\alpha\in A}x_\alpha$, prenons pour $\unsure{\forall|\mathrm{tt}}\alpha$ un foncteur $\unsure{\varphi|\phi}_\alpha:I_\alpha\unsure{\longrightarrow|\mathrel{-}}C$, avec $\operatorname{card}\mathrm{Fl}\,I_\alpha\leq\pi$, tel que $x_\alpha$ soit la $\unsure{\varinjlim|\underline{\lim}}$ dans $M$ de ce foncteur, on trouve que $x$ est $\varinjlim$ du foncteur $\varphi:I\longrightarrow C$ défini par les $\varphi_\alpha$, où $I=\coprod_\alpha I_\alpha$ (somme dans $\mathrm{Cat}$). Il est évident que $\operatorname{card}\mathrm{Fl}(I)\leq\pi\cdot\pi=\pi$, ce qui achève de prouver 2\textsuperscript{o}.

% element: text | foot, lines 18-22 | Prouvons 3 degrees
\uline{Prouvons} 3\textsuperscript{o}. On a les implications tautologiques
% element: deletion | foot, cancelled implication chain | 1 implies 1 prime implies 2
\deleted{$1)\unsure{\Leftrightarrow|\Rightarrow}1')\Rightarrow2$}
\inserted{$1\Rightarrow1'\unsure{)|,}$ et $\unsure{2'|1'}\Rightarrow1$} puisque $\mathrm{Filt}'{}^\pi$ est stable par sommes de cardinal $\leq\pi$. Donc $1\Leftrightarrow2'$, et implique que pour $\forall x$ dans $M$,
